\section*{章末练习}

\begin{longtable}{L{0.08\linewidth}L{0.32\linewidth}L{0.5\linewidth}}
\toprule
题号 & 类型 & 题目 \\
\midrule
1 & 补码 & 写出 8-bit 补码中 $-37$ 的 bit 模式，并计算它与 20 相加的结果。 \\
2 & 溢出 & 8-bit 补码中 $100+60$ 的 bit 结果是什么？进位标志与有符号溢出标志分别如何判断？ \\
3 & 加法器 & 比较 ripple-carry 与前缀/超前进位结构的延迟和布线代价。 \\
4 & 移位 & 算术右移与逻辑右移对负数有何不同？说明符号扩展来自哪里。 \\
5 & 乘法 & 用移位加法解释无符号乘法器需要哪些寄存器与控制状态；为何乘积通常需要双倍位宽？ \\
6 & 浮点边界 & 为什么浮点加法要先对阶？大数与很小的数相加时，后者可能怎样消失？ \\
\bottomrule
\end{longtable}

\section*{参考解答与要点}

\begin{longtable}{L{0.08\linewidth}L{0.32\linewidth}L{0.5\linewidth}}
\toprule
题号 & 类型 & 解答要点 \\
\midrule
1 & 补码 & 37 为 \code{00100101}，取反加一得 \code{11011011}；加 20 的 \code{00010100} 得 \code{11101111}，解释为 $-17$。 \\
2 & 溢出 & bit 结果为 160 的低 8 bit，即 \code{10100000}；无符号有向第 9 位的进位可独立判断，有符号数两个正数得到负号，故发生有符号溢出。 \\
3 & 加法器 & ripple-carry 面积与布线规则，但进位逐位传播；前缀/超前结构用更多生成传播逻辑和跨位连线缩短关键路径。 \\
4 & 移位 & 逻辑右移补 0，算术右移复制最高符号位，使补码负数在不考虑舍入时近似除以二。 \\
5 & 乘法 & 保存被乘数、乘数、累加部分积与迭代计数；每轮按乘数位选择加法并移位。两个 $n$ bit 无符号数最大乘积需要最多 $2n$ bit。 \\
6 & 浮点边界 & 两数尾数须对应同一指数再相加；指数差很大时小数尾数右移超过有效精度，舍入后可能对结果没有贡献。 \\
\bottomrule
\end{longtable}
